\chapter{Confined Spaces Risk Assessment}
\label{chap:Confined Spaces Risk Assessment}

%---------------------------- SECTION ----------------------------
\section{Why this assessment matters in construction}

\paragraph{Confined spaces are among the most lethal working environments encountered on construction sites, and their danger is characterised by a degree of invisibility that makes them uniquely difficult to manage through informal supervision and site culture. A manhole that has been entered safely a hundred times, a duct that feels unremarkable, a tank that appears clean and empty --- each can, without any warning, contain an atmosphere capable of killing an entrant within seconds of descent. The Confined Spaces Regulations 1997 exist precisely because the hazards of confined spaces are not reliably detectable by unaided human perception: an oxygen-deficient atmosphere is odourless, colourless, and produces no warning before the entrant collapses. The HSE's statistics on confined space fatalities consistently reflect a pattern in which the rescuer dies alongside the original victim --- a phenomenon driven by the instinct to assist a fallen colleague without pausing to don respiratory protection or to consider why the colleague fell. Multi-fatality confined space incidents, in which two, three, or four people die in the same space during a single event, are not rare occurrences: they are a recurrent pattern documented in HSE investigations over decades, and every such investigation reaches the same conclusion --- the deaths were preventable by the application of the controls specified in the Confined Spaces Regulations 1997 and the guidance in HSG258.}

\paragraph{The definition of a confined space under the Confined Spaces Regulations 1997 is deliberately broad: any place which is substantially enclosed (whether or not at an atmospheric pressure) and where there is a reasonably foreseeable specified risk. The specified risks are defined in Regulation~1(1) and include: a serious injury to any person at work there arising from a fire or explosion; loss of consciousness arising from an increase in body temperature; loss of consciousness or asphyxiation arising from gas, fume, vapour, or the lack of oxygen; drowning arising from an increase in the level of a liquid; asphyxiation arising from a free-flowing solid; and loss of consciousness or asphyxiation in any of the circumstances described. This definition encompasses a far wider range of locations on a construction site than the instinctive picture of a deep underground space. Shallow manholes, drainage ducts, roof voids, basement plant rooms during commissioning, storage tanks, excavations deeper than approximately 1.2 metres in certain soil and ground conditions, cable ducts, building service risers, and even large vessels and vessels partially submerged below grade all fall within the regulatory definition if a specified risk is reasonably foreseeable. The assessment process must not begin with a pre-formed mental image of what a confined space looks like and then attempt to match the site conditions to that image --- it must begin with a systematic hazard identification process applied to every substantially enclosed space on the site.}

\paragraph{The consequences of a confined space incident are characterised by their speed of onset and their irreversibility. An oxygen-deficient atmosphere at 16\% oxygen will cause loss of consciousness in an entrant within approximately two minutes; at 6\% oxygen, unconsciousness is essentially instantaneous. Carbon monoxide (CO) at 1,200 ppm will cause loss of consciousness and death within one hour; CO at concentrations produced by petrol-engine plant operating in a confined space can exceed 5,000 ppm and deliver a fatal dose in minutes. Hydrogen sulphide (H\textsubscript{2}S), which is produced in sewage and decomposing organic material, has a distinctive unpleasant odour at low concentrations but causes rapid olfactory fatigue: an entrant may initially detect the odour of H\textsubscript{2}S and believe they have adequate warning, but within minutes the olfactory nerve is overcome and the entrant loses the ability to detect the gas at any concentration. High concentrations of H\textsubscript{2}S (above approximately 700 ppm) produce immediate collapse without olfactory warning. These characteristics make multi-gas monitoring before and during entry an absolute requirement, not a discretionary enhancement: the human senses are not a reliable substitute for calibrated atmospheric monitoring instruments.}

\paragraph{The construction industry encounters confined spaces in virtually every project type and at every phase of construction. Groundworks and drainage phases involve entry into manholes, sewers, and inspection chambers, often during live drainage connections. Structural phases involving deep basements, underground car parks, and service tunnels create large-scale confined environments. Mechanical and electrical phases involve entry into plant rooms, ductwork, riser shafts, and storage vessels. Refurbishment and demolition phases involve entry into tanks, voids, and underground structures whose history of use may be unknown, and whose atmospheric condition may have been affected by contaminated ground or residual chemical deposits. The spread of confined space risk across all project phases means that a confined space risk assessment is not a document required only for specialist contractors engaged on drainage or utilities work: it is a document required on essentially every construction project and must be treated as a primary document within the Construction Phase Plan.}

\barquote{``In a confined space, the difference between a safe entry and a fatality is measured not in minutes but in seconds. The standby person, the rescue plan, and the atmospheric test are not optional enhancements --- they are the margin between life and death, and they must be in place before anyone descends.''}

\example{Confined Space Fatality Scenario}{A drainage subcontractor is connecting a new surface water drainage system to an existing combined sewer network on a housing development. A manhole on the existing combined sewer, approximately 2.4 metres deep, requires internal inspection and connection of the new branch pipe. The manhole had last been entered two years previously during the pre-construction survey and had been recorded as clear at that time. The drainage operative descends the manhole on a Monday morning without conducting any pre-entry atmospheric testing, relying on the previous survey record and his own assessment that the manhole ``looks clean.'' At the base of the manhole, he collapses without warning. His colleague, observing from the surface, descends to assist him and also collapses. A passing plant operative sees both workers in the manhole and descends to assist, and is rendered unconscious within approximately 30 seconds. Emergency services attend and conduct a non-entry rescue using a tripod and winch, recovering all three workers. One drainage operative dies; the other two workers suffer hypoxic brain injury. Atmospheric testing of the manhole at the time of rescue records oxygen at 10.8\% and carbon dioxide at approximately 4\%, consistent with oxygen displacement by biogenic CO\textsubscript{2} from decomposing organic material in the sewer network. There was no confined space entry permit, no standby person, no rescue equipment at the worksite, and no pre-entry atmospheric test. The incident is entirely consistent with the pattern described in HSG258 and demonstrates every element of a preventable multi-fatality confined space event.}

%---------------------------- SECTION ----------------------------
\section{Legal and professional framework}

\paragraph{The legal framework for confined space working on construction sites is anchored by the Confined Spaces Regulations 1997 (CSR 1997), which apply to all workplaces including construction sites and impose specific duties on employers, the self-employed, and those in control of premises in relation to work in confined spaces. The CSR 1997 impose a clear hierarchy of duty: first, avoid entry to a confined space if it is reasonably practicable to do so by performing the work from outside the space; second, if entry cannot be avoided, follow a safe system of work for entry; third, ensure that adequate arrangements for rescuing entrants are in place before entry commences. The regulations are supported by the Approved Code of Practice and guidance document L101 (Safe Work in Confined Spaces, third edition, 2014), which provides authoritative interpretation of the regulatory requirements and detailed guidance on atmospheric monitoring, rescue arrangements, and safe systems of work. The HSE's Guidance on Regulations document HSG258 (Controlling Risks in Confined Spaces) provides additional operational guidance for specific confined space environments including sewers, tanks, and vessels. The Construction (Design and Management) Regulations 2015 require that confined space risks are identified in the pre-construction information, addressed in the Construction Phase Plan, and controlled by competent contractors whose method statements demonstrate compliance with the CSR 1997 and L101 framework.}

\paragraph{Key frameworks relevant to this chapter include: Confined Spaces Regulations 1997; Health and Safety at Work etc.\ Act 1974; Management of Health and Safety at Work Regulations 1999; Construction (Design and Management) Regulations 2015; L101 (Safe Work in Confined Spaces); HSG258 (Controlling Risks in Confined Spaces); RIDDOR 2013; BS~EN~250 (respiratory protective equipment); BS~EN~137 (self-contained open-circuit compressed air breathing apparatus); BS~EN~14593 (compressed air line breathing apparatus); and City \& Guilds 6150-03 confined space training standard.}

\begin{itemize}
  \bitem{Confined Spaces Regulations 1997 (CSR 1997)}{The primary statutory instrument governing work in confined spaces in Great Britain. Regulation~4 requires that no person at work shall enter a confined space to carry out work for any purpose unless it is not reasonably practicable to achieve that purpose without such entry. Regulation~5 requires that, where entry is necessary, the employer must ensure a suitable safe system of work is followed, and that suitable rescue arrangements are made before entry commences. Regulation~6 requires that rescue arrangements include equipment capable of recovering an entrant from the confined space without requiring the rescuer to enter the space, where reasonably practicable. The Approved Code of Practice (L101) accompanying the regulations sets out the specific requirements for atmospheric testing, standby persons, communications, and rescue equipment that constitute compliance.}
  \bitem{L101 --- Safe Work in Confined Spaces (HSE, third edition 2014)}{The Approved Code of Practice and guidance issued under the CSR 1997. L101 provides authoritative and detailed guidance on: identifying confined spaces and specifying the associated risks; establishing safe systems of work including confined space entry permits; pre-entry atmospheric testing requirements and monitoring thresholds; selection and use of respiratory protective equipment (RPE) for confined spaces; standby person duties and competence; emergency and rescue procedures; and training standards for confined space workers. Compliance with L101 creates a presumption of compliance with the CSR 1997. L101 specifically states that non-entry rescue should always be the preferred option and that entry rescue should only be conducted by trained rescue teams.}
  \bitem{HSG258 --- Controlling Risks in Confined Spaces (HSE)}{HSE guidance document providing operational detail supplementary to L101, with specific treatment of confined space working in sewers and drains, tanks and vessels, and excavations. HSG258 provides guidance on atmospheric monitoring in sewers (addressing the specific hazards of CO, H\textsubscript{2}S, oxygen deficiency, and flammable gases in sewer environments); the use of tripod and winch systems for non-entry rescue; communication systems for confined space entry; and the management of confined space working by contractors. The guidance also addresses specific scenarios such as entry into live sewers, entry into tanks containing hazardous residues, and entry in areas where ground contamination may affect the atmospheric quality of the confined space.}
  \bitem{BS EN 60079 series and ATEX 153 --- Explosive Atmospheres}{Where a confined space may contain a flammable or explosive atmosphere (arising from flammable gases, vapours, or dusts), ATEX zone classification under the Equipment and Protective Systems for Use in Potentially Explosive Atmospheres Regulations 2016 applies. Monitoring equipment, lighting, communications equipment, and any other electrical or electronic equipment used within or near the space must be rated for use in the applicable ATEX zone. Gas monitoring instruments used in potentially explosive atmospheres must themselves be intrinsically safe and carry the appropriate ATEX certification marking.}
  \bitem{Personal Protective Equipment at Work Regulations 2022 (PPER 2022)}{Requires that appropriate PPE, including respiratory protective equipment suitable for the hazards present in the confined space, is provided to workers at no cost and that its selection, use, maintenance, and storage are managed in accordance with a written risk assessment. For confined spaces, RPE selection must address the specific atmospheric hazards identified by pre-entry testing and must provide adequate protection against the worst credible atmospheric condition, not merely the expected normal condition.}
  \bitem{Provision and Use of Work Equipment Regulations 1998 (PUWER 1998)}{Applies to all work equipment used in confined space operations, including multi-gas detectors, tripod and winch rescue systems, self-contained breathing apparatus (SCBA), airline breathing apparatus, and communication equipment. PUWER requires that equipment is suitable for its intended purpose, maintained in an efficient state, regularly inspected, and operated only by trained and competent persons. Multi-gas detectors must be calibrated and bump-tested before each use in accordance with manufacturers' instructions and the requirements of L101.}
  \bitem{RIDDOR 2013}{Requires the immediate notification to the relevant enforcing authority of any fatality or specified injury arising from a confined space incident. Confined space incidents resulting in loss of consciousness, hospitalisation, or death must be reported. The RIDDOR reporting obligation applies to the person in control of the premises, which on a construction site will typically be the principal contractor. Any confined space incident must also trigger an immediate investigation and review of the confined space entry permit system and risk assessment.}
  \bitem{City \& Guilds 6150-03 and IATP-equivalent Confined Space Training}{The Confined Spaces Regulations 1997 and L101 require that persons who enter confined spaces and those who act as standby persons are suitably trained and competent. Industry-recognised confined space training qualifications include City \& Guilds 6150-03 (Award in Confined Space Operations at Low Risk and Medium Risk), with progression to High Risk and Emergency Response and Rescue qualifications for higher-risk environments. Training must be specific to the level of risk and the type of confined space being entered, and must include practical assessment in realistic simulated confined space conditions.}
\end{itemize}

%---------------------------- SECTION ----------------------------
\section{Conducting the assessment using the five-step method}

\subsection{Step 1 --- Identify Hazards}
\paragraph{Hazard identification for confined space risk assessment requires both desk-based research and physical site survey. The desk-based element should draw on the pre-construction information provided under CDM 2015, which must include details of underground drainage and service infrastructure, ground contamination records, former use of the site (which may indicate the presence of tanks, voids, or underground structures not shown on current plans), and any utility-supplied records of below-ground assets. The physical survey should identify and record every substantially enclosed space on the site, however small or apparently innocuous, and assess whether a specified risk under the CSR 1997 is reasonably foreseeable in that space. Hazard categories to be addressed in the identification stage include: toxic atmosphere, arising from decomposition of organic matter in sewers and drains (H\textsubscript{2}S, CO\textsubscript{2}, CO), from ground contamination (volatile organic compounds, methane), from residual chemicals in tanks, from combustion products of plant operating near or within the space, or from chemical reactions within the space; oxygen deficiency (O\textsubscript{2} below 19.5\%), arising from biological consumption, displacement by inert gases, absorption by rusting metal surfaces, or microbial activity in organic material; oxygen enrichment (O\textsubscript{2} above 23.5\%), arising from leaking oxygen cylinders or oxidant pipework, which dramatically increases the flammability of clothing and equipment and can cause spontaneous combustion of materials; flammable or explosive atmosphere, arising from methane in sewer gas, flammable vapours from chemical residues, LPG from leaking equipment, or natural gas ingress from adjacent pipework; liquid ingress, from flooding of the space by surface water, groundwater, or upstream flow in a live sewer; excessive heat, from inadequate ventilation in summer conditions, heat-generating equipment within the space, or proximity to heat sources in adjacent building fabric; free-flowing solids, in spaces adjacent to grain storage, sand storage, or areas where particulate material may flow into the space through a drainage connection or vent.}

\subsection{Step 2 --- Decide Who or What Is Affected}
\paragraph{The population at risk from confined space hazards on a construction site includes: drainage operatives carrying out connection, inspection, and maintenance work in manholes, sewers, inspection chambers, and pump sumps; mechanical and electrical operatives working in underground plant rooms, service ducts, cable vaults, and riser shafts; groundworkers entering deep excavations in ground conditions that may give rise to specified risks (such as made ground with methane potential or ground adjacent to sewers); specialist subcontractors engaged for specific confined space operations including sewer rehabilitation, tank cleaning, and duct installation; and the standby person, who is at risk from the instinct to enter the space in the event of an emergency if rescue procedures and equipment are not clearly specified and practised. Persons conducting rescue, including site first-aiders, emergency services personnel, and untrained bystanders, are at very high risk if they attempt an unprotected entry rescue. The assessment must explicitly address the risk to would-be rescuers and must establish the non-entry rescue system and its location as an absolute prerequisite of entry. Visiting professionals including utilities engineers attending for service connections, and client representatives requiring access to inspect underground drainage structures, must not be permitted to enter a confined space without being subject to the same assessment and permit controls as the primary workforce.}

\subsection{Step 3 --- Evaluate Likelihood and Consequence}
\paragraph{The consequence of a toxic or oxygen-deficient atmosphere event in a confined space is loss of consciousness and death, and this consequence is not theoretical --- it is the documented outcome of the majority of confined space incidents recorded by the HSE. The inherent risk of entry into an untested confined space must therefore be rated at the highest level on any risk matrix before controls are applied: Very High likelihood of fatal consequences if entry is made without prior atmospheric testing and a safe system of work. The likelihood evaluation must resist any reduction based on the history of previous uneventful entries, because the confined space hazard is not attenuated by frequency of exposure --- an atmosphere that has been safe on every previous entry can become lethal during the interval between entries if a change in the drainage system, ground conditions, or ambient temperature has altered the atmospheric composition. The pre-entry atmospheric test is the only mechanism by which the likelihood of a toxic or deficient atmosphere is reduced from the inherent worst-case assumption: without the test result, the likelihood must be treated as maximum. Specific likelihood factors that increase the risk profile include: spaces connected to live sewers (high likelihood of H\textsubscript{2}S and CO\textsubscript{2}); spaces in areas of known ground contamination (potential for methane or volatile organic compounds); spaces that have been closed and unventilated for extended periods (risk of oxygen depletion by absorption or biological activity); spaces in hot weather (increased rate of biological activity and vapour generation); and spaces adjacent to or below areas where petrol or diesel plant is operating (risk of CO ingress from exhaust gases).}

\subsection{Step 4 --- Select and Implement Controls}
\paragraph{The primary control hierarchy for confined spaces begins with avoidance: the first question the risk assessment must answer is whether the work can be carried out without entering the space. Remote inspection using CCTV survey cameras eliminates the need for manned entry into drainage structures for inspection purposes; jet-wash and robotic rehabilitation systems enable drainage cleaning without entry; remote actuators and motorised valves can be operated from outside plant rooms; extended tools can be used for certain fixing operations in duct access. Where avoidance is genuinely impracticable, the safe system of work for entry must incorporate the following mandatory elements as specified in L101 and CSR 1997. Pre-entry atmospheric testing using a calibrated multi-gas monitor capable of measuring oxygen percentage, flammable gases as a percentage of the lower explosive limit (LEL), carbon monoxide in ppm, and hydrogen sulphide in ppm must be conducted at the top, middle, and bottom of the confined space before any person enters. The established safe working thresholds are: oxygen 19.5--23.5\%; flammable gases below 10\% LEL (entry prohibited above this level); CO below 20 ppm STEL (short-term exposure limit); H\textsubscript{2}S below 5 ppm TWA. Continuous atmospheric monitoring by the entrant during the period of occupation must be maintained throughout the entry. The confined space entry permit must be completed and signed by the appointed supervisor before entry commences, specifying the space to be entered, the work to be carried out, the atmospheric test results, the RPE selected, the communication system, the rescue equipment in place, and the name of the standby person. The standby person must be present at the access point throughout the entry, must be in continuous communication with the entrant, and must not enter the space under any circumstances. The rescue arrangement must be established before entry, with the tripod and winch (or equivalent non-entry rescue device) rigged and ready. RPE for confined spaces must be selected to match the atmospheric hazard: self-contained breathing apparatus (SCBA) or airline breathing apparatus (continuous flow) is required where the atmosphere is immediately dangerous to life and health (IDLH) or where the oxygen content is unknown or deficient; supplied-air respirators and airline systems are appropriate for extended work periods where a fixed air supply can be provided; filtering face pieces are not appropriate where oxygen deficiency or toxic gas hazard is present.}

\subsection{Step 5 --- Record and Review}
\paragraph{The confined space entry permit is the central recording document for each individual entry operation and must be treated as a live document throughout the period of entry. The permit must record: the date and time of entry; the name of the entrant and the standby person; the name of the authorising supervisor; the atmospheric test results before entry, with instrument identification number and calibration date; the RPE type and identification number issued; the communication method established; the rescue equipment in place and its location; any specific hazards identified for the particular entry; and the time of exit and confirmation that the permit has been closed out. Completed permits must be retained as part of the Construction Phase Plan records and are disclosable in the event of an incident. The confined space risk assessment itself must be reviewed at the start of each day on which confined space entries are planned, taking account of any changes to the drainage system, ground conditions, adjacent works, weather, or any other factor that may have altered the atmospheric condition of the space since the previous assessment. Any incident, near miss, atmospheric alarm actuation, or operational deviation from the permit must trigger an immediate formal review of the risk assessment before the next entry is authorised. The review must include re-calibration and bump-testing of the atmospheric monitoring equipment.}

\example{Five-Step Application}{A civil engineering contractor is constructing a new surface water attenuation tank system for a logistics park. The scheme involves the installation of large-diameter pre-cast concrete culvert sections, 2.4 metres internal diameter, partially buried in ground with a known history of landfill on part of the site. The contractor must enter the culvert sections during installation to make rubber ring joint connections. \textbf{Step~1 (Hazards):} Oxygen deficiency due to absorption by the concrete; methane and CO\textsubscript{2} from adjacent landfill gas migration; potential H\textsubscript{2}S from organic material in the landfill deposit; limited natural ventilation in completed sections; liquid ingress from groundwater dewatering pump failure. \textbf{Step~2 (Who Affected):} Jointing operatives; banksman at access point; groundworks foreman; any operative passing within proximity of the open ends during other groundworks activities. \textbf{Step~3 (Evaluation):} Inherent risk rated Very High for toxic atmosphere and oxygen deficiency given landfill gas potential; Very High for liquid ingress without dewatering. \textbf{Step~4 (Controls):} Pre-entry four-gas atmospheric test conducted at all three points (entry, mid, end) before each entry, with entry prohibited if O\textsubscript{2} below 19.5\%, flammable gases above 10\% LEL, CO above 20 ppm, or H\textsubscript{2}S above 5 ppm; continuous monitoring by entrant using personal four-gas clip-on monitor throughout occupation; forced ventilation using a venturi fan providing fresh air flow throughout; confined space entry permit raised for each entry signed by site manager; standby person designated and briefed, stationed at access point; tripod and winch rigged at entry point before each entry; entrant to wear full body harness connected to the winch line; SCBA on standby with standby person trained in its use; dewatering pump monitored continuously during occupation of any section. \textbf{Step~5 (Record):} Permit raised, signed, and closed for each entry; atmospheric test results recorded on permit with instrument serial number; permits retained in Construction Phase Plan; risk assessment reviewed daily before first entry.}

%---------------------------- SECTION ----------------------------
\section{Applying the hierarchy of controls}

\paragraph{The hierarchy of controls for confined spaces is given statutory expression in the CSR 1997 Regulation~4, which establishes avoidance as the primary duty and safe systems of work as the secondary duty. The hierarchy must be applied rigorously: entry into a confined space must never be treated as the default option simply because it is the most familiar approach. The application of each level of the hierarchy must be documented in the risk assessment and in the method statement, so that the decision to enter the space is demonstrably a last resort rather than an unreflective default.}

\begin{itemize}
  \bitem{Elimination}{Avoid entry to the confined space entirely where it is reasonably practicable to do so. CCTV drainage surveys, robot-operated cleaning and lining systems, remote-operated valve actuators, and extended-reach tooling can perform many functions that previously required manned entry. Where a drainage inspection manhole requires an internal visual check, a handheld CCTV camera lowered from the surface eliminates the entry requirement. Designers specifying maintenance arrangements for below-ground assets should consider the confined space avoidance duty under CSR 1997 Regulation~4 during the CDM design process, and design out the need for routine manned entry wherever practicable through the provision of adequate access covers, adequate remote inspection capability, and adequate service routing to permit above-ground maintenance.}
  \bitem{Substitution}{Where entry cannot be entirely eliminated, consider whether the work can be carried out by a substitute method that reduces the duration of entry, the number of entries required, or the severity of the atmospheric hazard. Substituting mechanical jointing methods for labour-intensive manual jointing inside a culvert reduces the time an operative must spend inside the confined space; pre-assembling drainage modules above ground before lowering as a unit reduces the number of individual entry events; providing forced ventilation before entry and throughout the operation substitutes the contaminated internal atmosphere with a supply of fresh air, reducing the atmospheric hazard from the inherent worst case to a controlled managed level.}
  \bitem{Engineering controls}{Where entry must proceed, engineering controls include: mechanical forced ventilation to dilute and remove atmospheric contaminants and maintain oxygen levels within the safe range throughout the entry; atmospheric monitoring by fixed-point instruments or by personal clip-on monitors worn by entrants, with audible and visual alarms at the pre-set threshold concentrations; tripod and winch non-entry rescue systems rigged at the access point before entry commences and capable of recovering an incapacitated entrant; gas purging of tanks and vessels containing known contaminated atmospheres before entry is attempted; explosion-proof lighting and intrinsically safe electrical equipment in spaces where flammable atmospheres may be present; and locked-out and tagged-out isolation of all energy sources (electrical, hydraulic, pneumatic, thermal, and gravity) feeding into the confined space before entry.}
  \bitem{Administrative controls}{The confined space entry permit system is the primary administrative control, establishing the authority to enter, the conditions for entry, and the exit requirements as a formally controlled process. Additional administrative controls include: training all persons involved in confined space operations (entrants, standby persons, and rescue teams) to the standard required by L101 and City \& Guilds 6150-03 or equivalent; designating a named supervisor responsible for authorising each entry; establishing a check-in and check-out system with an agreed maximum entry duration before the standby person initiates the rescue protocol; specifying the communication system (two-way radio, signal rope, or direct verbal contact) and the alarm signals; and briefing all site personnel not involved in confined space operations on the location of active confined spaces and on the absolute prohibition on unauthorised entry.}
  \bitem{PPE}{Respiratory protective equipment for confined space entry must be selected against the specific atmospheric hazard identified by pre-entry testing. Self-contained breathing apparatus (SCBA) to BS~EN~137 provides a fully independent air supply unaffected by the ambient atmosphere and is required where the atmospheric condition is unknown, where IDLH concentrations of toxic gases have been measured, or where oxygen deficiency is present or possible. Compressed air line breathing apparatus (airline BA) to BS~EN~14593 provides a continuous supply of compressed air via a hose from an external supply and is suitable for extended work periods where the supply hose length is compatible with the working geometry. Gas-tight chemical protective suits may be required where the internal surfaces of tanks contain toxic or corrosive chemical residues. Entrants must wear a full-body harness connected to the non-entry rescue system at all times during occupancy of the confined space. Standard filtering face pieces and half-mask respirators are not appropriate for confined space entry where oxygen deficiency is possible.}
\end{itemize}

%---------------------------- SECTION ----------------------------
\section{Cadence of assessment and review triggers}

\paragraph{Confined space risk assessments must be treated as dynamic documents subject to review before every entry event and after any change that may affect the atmospheric or physical condition of the space. The confined space entry permit serves as the operational review mechanism for each individual entry, but the underlying risk assessment must also be formally reviewed at appropriate intervals and following the specific trigger events described below. The cadence of atmospheric testing must reflect the known rate of change of conditions in the specific space: a sewer manhole adjacent to a live combined sewer can experience rapid and unpredictable changes in atmospheric composition as flow conditions in the sewer change, and continuous monitoring throughout entry is non-negotiable in such environments.}

\begin{itemize}
  \bitem{Before every entry}{The confined space entry permit must be raised for each individual entry event. This necessarily includes a review of the risk assessment by the authorising supervisor, confirmation that the atmospheric test has been conducted and results are within safe limits, confirmation that rescue equipment is rigged and standby person is briefed and in position, and confirmation that no changes to the drainage system, ground conditions, or adjacent works have altered the risk profile since the last entry.}
  \bitem{After any change to drainage or service connectivity}{Any upstream drainage connection, valve operation, storm flow event, or change to the sewerage network in the catchment served by the space may alter the composition of gases in the space within minutes. Sewer gas composition is highly variable and is influenced by water temperature, organic loading, turbulence, and ventilation. Any change to the connected infrastructure must trigger re-testing before the next entry, regardless of how recently the previous test was conducted.}
  \bitem{After adverse weather or temperature change}{High ambient temperatures increase the rate of biological gas production in organic material in sewers and excavations. Heavy rainfall can cause rapid rises in water level in sewers and sumps and can displace accumulated gas from sewer networks into manholes and chambers. Both scenarios require re-testing of the confined space atmosphere before entry following these weather events.}
  \bitem{Following any atmospheric alarm actuation}{Any actuation of an atmospheric monitor alarm during an ongoing entry must result in immediate evacuation of the space, assessment by the authorising supervisor, re-testing, and formal review of the risk assessment before re-entry is permitted. An alarm actuation is evidence that the atmospheric condition of the space has changed unexpectedly and that the original risk assessment basis may no longer be valid.}
  \bitem{Following any incident, near miss, or rescue event}{Any confined space incident, near miss, alarm actuation, or rescue operation must trigger a full formal review of the risk assessment, the permit system, the training of the personnel involved, the adequacy of the rescue equipment, and the supervision arrangements. No further entry into the space should be authorised until the review is complete and corrective actions are implemented and documented.}
\end{itemize}

%---------------------------- SECTION ----------------------------
\section{Common failure modes}

\begin{itemize}
  \bitem{Entry without pre-entry atmospheric testing}{The single most common causal factor in confined space fatalities is entry into the space without conducting any atmospheric test. The reasons given in investigations are consistent: the space appeared normal, previous entries had been uneventful, the test equipment was not immediately available, or the operative did not consider the space to be a confined space requiring a permit. Pre-entry atmospheric testing is not a formality --- it is the only mechanism by which the unknown atmospheric condition of the space is made known before a human enters it. The decision to enter without testing is the decision to accept an unknowable risk that has killed many times previously.}
  \bitem{Use of oxygen-only indicators rather than multi-gas monitors}{A significant number of confined space incidents have involved situations where operatives checked the oxygen level using a single-sensor oxygen indicator, found it within the normal range, and concluded the atmosphere was safe. An atmosphere can contain a normal oxygen level (approximately 20.9\%) while simultaneously containing lethal concentrations of CO, H\textsubscript{2}S, or other toxic gases. The requirement for multi-gas monitoring --- oxygen, flammable gases, CO, and H\textsubscript{2}S as a minimum --- is specifically addressed in L101 for this reason. Single-parameter monitors do not constitute adequate pre-entry testing.}
  \bitem{Failure to bump-test and calibrate monitoring equipment}{Multi-gas monitors are electrochemical instruments whose sensors degrade over time and may fail without obvious indication. Bump testing (exposing the monitor to a known concentration of test gas to verify sensor response) before each use is mandatory under L101 and the manufacturers' instructions. Monitors not bump-tested before use may give a safe reading even when the actual atmospheric concentration of a toxic gas is lethal. The monitoring equipment log, including calibration certificate dates and bump test records, must be available for inspection at the point of entry.}
  \bitem{Untrained rescue or bystander entry following an emergency}{The majority of confined space multi-fatality incidents involve one or more deaths among persons attempting to rescue the original victim without respiratory protection. The instinct to assist a fallen colleague is powerful and commendable, but entry into a confined space containing a toxic or oxygen-deficient atmosphere without SCBA is invariably fatal. The standby person must be explicitly briefed that their absolute duty is to remain outside the space, to initiate the non-entry rescue using the tripod and winch, and to contact the emergency services --- not to enter the space. This instruction must be given clearly, in writing, as part of the permit briefing, and must be reinforced in training.}
  \bitem{Inadequate or incorrectly rigged rescue equipment}{The non-entry rescue system (tripod, winch, and rescue line attached to the entrant's harness) must be rigged and ready before the first entrant descends. A tripod and winch stored in a van 50 metres away from the access point does not constitute an adequate rescue arrangement. The system must be capable of recovering an incapacitated entrant from the maximum depth and configuration of the space, and the standby person must be trained and practised in its use. Regularly recurring investigation findings include: rescue winch not attached to entrant's harness; entrant not wearing a harness; tripod not rated for the recovery load; and standby person not aware of how to operate the winch mechanism.}
  \bitem{Reliance on natural ventilation in deep or complex spaces}{Natural ventilation of confined spaces, while potentially adequate in simple shallow open-top spaces with good air circulation, is frequently insufficient in deep manholes, long drainage runs, enclosed tanks, and any space with restricted air flow geometry. The assumption that a space is adequately ventilated because it ``has an opening'' is not acceptable as the basis for omitting forced ventilation. Forced mechanical ventilation using a venturi or axial fan directing fresh air into the base of the space must be specified for any space where the depth, configuration, or atmospheric hazard assessment indicates that natural air change rates may not maintain safe oxygen levels and dilute toxic gases throughout the entry period.}
\end{itemize}

\example{Failure Pattern}{An investigation into a confined space incident at a new-build commercial development identified the following failure chain. A drainage gang was carrying out internal joint inspection and sealing inside a newly installed large-diameter concrete sewer pipe in a 3.2-metre-deep trench. The sewer had been connected to the existing public sewer at one end the previous day, creating a live connection to the live sewer network. No confined space assessment was conducted because the sewer pipe was perceived as a ``new'' asset rather than a confined space. No atmospheric test was conducted. No standby person was designated. No entry permit was raised. One operative entered the pipe at 7.45\,am without any pre-entry preparation. He collapsed approximately 12 metres from the entry point. His colleague entered immediately upon noticing his colleague had not returned and also collapsed. Both were recovered by emergency services using a recovery line that had not been pre-rigged but was improvised by the fire service on attendance. One operative died; one sustained permanent neurological injury. Atmospheric testing during recovery recorded O\textsubscript{2} at 13.2\%, consistent with oxygen displacement by CO\textsubscript{2} and H\textsubscript{2}S from the connected live sewer. The investigation identified: no confined space identification process; no awareness that the live connection had fundamentally altered the atmospheric status of the new pipe; no standby person; no rescue equipment; and no training of any of the gang members in confined space requirements. The principal contractor had not requested a confined space method statement from the drainage subcontractor and had not verified compliance with CSR 1997.}

%---------------------------- CHAPTER SUMMARY ----------------------------
\newpage
\section{Chapter Summary}

\begin{table}[H]
\centering
\begin{tabularx}{\textwidth}{>{\raggedright\arraybackslash}p{3.2cm}X}
\toprule
\textbf{Assessment Element} & \textbf{Operational Requirement} \\
\midrule
Legal/Professional Basis & Confined Spaces Regulations 1997; L101 (Safe Work in Confined Spaces, 3rd ed.\ 2014); HSG258 (Controlling Risks in Confined Spaces); Health and Safety at Work etc.\ Act 1974; RIDDOR 2013; CDM 2015; City \& Guilds 6150-03 or equivalent confined space training. \\
Method & Apply the full five-step process and document inherent/residual risk. Raise a confined space entry permit for each individual entry; conduct calibrated multi-gas atmospheric testing before every entry and maintain continuous monitoring throughout. \\
Controls & Avoidance of entry as primary control; forced mechanical ventilation; multi-gas pre-entry and continuous monitoring (O\textsubscript{2} 19.5--23.5\%, flammable gases $<$10\% LEL, CO $<$20 ppm, H\textsubscript{2}S $<$5 ppm); confined space entry permit; trained standby person stationed at access point throughout; non-entry rescue system (tripod and winch) rigged before entry; SCBA or airline BA for IDLH environments. \\
Cadence & Review before every entry via the permit process; re-test after any drainage system change, adverse weather, temperature change, or alarm actuation; full risk assessment review following any incident or near miss. \\
Assurance & Use audit, incident learning, and governance reporting to verify effectiveness. \\
\bottomrule
\end{tabularx}
\end{table}

\begin{itemize}
  \bitem{Key takeaway 1}{No person shall enter a confined space unless entry cannot be avoided --- the first control question is always whether the work can be performed from outside the space, and this question must be answered with evidence, not assumption.}
  \bitem{Key takeaway 2}{Pre-entry atmospheric testing using a calibrated multi-gas monitor (O\textsubscript{2}, flammable gases, CO, H\textsubscript{2}S) is mandatory before every entry and must be conducted at multiple levels within the space; no entry is to be made without a recorded test result within safe limits.}
  \bitem{Key takeaway 3}{The standby person must remain outside the confined space at all times and must not enter under any circumstances --- non-entry rescue using the rigged tripod and winch is the only acceptable first-response action; entry by an unprotected rescuer is the mechanism by which single-fatality incidents become multi-fatality incidents.}
  \bitem{Key takeaway 4}{All persons carrying out confined space entry or standby duties must be trained and certificated to the standard appropriate to the level of risk, as specified in L101 and the CSR 1997 Approved Code of Practice --- familiarity with the task is not an acceptable substitute for formal confined space training.}
  \bitem{Key takeaway 5}{The confined space entry permit must be raised, completed, and signed by the authorising supervisor before every entry, and must be closed out and retained on completion; a generic confined space risk assessment without a site-specific permit for each entry does not constitute compliance with the Confined Spaces Regulations 1997.}
\end{itemize}

\barquote{In Chapter~\ref{chap:Pressure Systems and Compressed Air Risk Assessment}, we examine the risks associated with pressure systems, compressed air equipment, and gas cylinders on construction sites, and the requirements of the Pressure Systems Safety Regulations 2000.}
